% ---------------------------------------------------------------------------
% Demo guide for the Autonomous Product Intelligence Factory (V2).
%
% Written to be followed by somebody who has never seen the application, and
% who has to present it to a CIO. Every command and every button name in here
% was checked against the running application rather than written from memory.
%
% Compile with any LaTeX engine:  pdflatex demo-guide.tex   (run it twice, so
% the table of contents fills in). Or paste the whole file into overleaf.com,
% which needs nothing installed.
% ---------------------------------------------------------------------------

\documentclass[11pt,a4paper]{article}

\usepackage[T1]{fontenc}
\usepackage[utf8]{inputenc}
\usepackage[margin=2.4cm]{geometry}
\usepackage{parskip}
\usepackage{xcolor}
\usepackage{graphicx}
\usepackage{enumitem}
\usepackage{tcolorbox}
\usepackage{fancyhdr}
\usepackage{titlesec}
\usepackage{longtable}
\usepackage{booktabs}
\usepackage[hidelinks]{hyperref}

\tcbuselibrary{skins,breakable}

% --- colours ---------------------------------------------------------------
\definecolor{ink}{HTML}{1B1F2A}
\definecolor{accent}{HTML}{2D4B8E}
\definecolor{quiet}{HTML}{6B7280}
\definecolor{shell}{HTML}{F3F4F6}
\definecolor{warnbg}{HTML}{FEF6E7}
\definecolor{warnline}{HTML}{C9821A}
\definecolor{saybg}{HTML}{EEF2FB}
\definecolor{okline}{HTML}{2F7D5C}

\color{ink}

% --- headings --------------------------------------------------------------
\titleformat{\section}{\Large\bfseries\color{accent}}{\thesection}{0.7em}{}
\titleformat{\subsection}{\large\bfseries}{\thesubsection}{0.6em}{}
\titlespacing*{\section}{0pt}{1.6em}{0.6em}

\pagestyle{fancy}
\fancyhf{}
\lhead{\small\color{quiet}Product Intelligence Factory --- Demo Guide}
\rhead{\small\color{quiet}\thepage}
\renewcommand{\headrulewidth}{0.4pt}

% --- reusable boxes --------------------------------------------------------

% Something to type into a terminal.
\newtcolorbox{cmd}{
  breakable, colback=shell, colframe=shell, boxrule=0pt,
  left=8pt, right=8pt, top=6pt, bottom=6pt, arc=2pt,
  fontupper=\ttfamily\small
}

% Words to say out loud.
\newtcolorbox{say}[1][]{
  breakable, colback=saybg, colframe=accent, boxrule=0pt,
  leftrule=3pt, left=8pt, right=8pt, top=6pt, bottom=6pt, arc=0pt,
  title=#1, coltitle=accent, fonttitle=\bfseries\small,
  attach title to upper=\par
}

% Something that will go wrong if ignored.
\newtcolorbox{careful}[1][Careful]{
  breakable, colback=warnbg, colframe=warnline, boxrule=0pt,
  leftrule=3pt, left=8pt, right=8pt, top=6pt, bottom=6pt, arc=0pt,
  title=#1, coltitle=warnline, fonttitle=\bfseries\small,
  attach title to upper=\par
}

% A fact worth knowing but not needed to run the demo.
\newtcolorbox{aside}[1][Background]{
  breakable, colback=white, colframe=quiet, boxrule=0.4pt,
  left=8pt, right=8pt, top=6pt, bottom=6pt, arc=2pt,
  title=#1, coltitle=quiet, fonttitle=\bfseries\small,
  attach title to upper=\par
}

\newcommand{\ui}[1]{\textbf{#1}}          % a thing on screen
\newcommand{\key}[1]{\textsf{\small #1}}  % a keyboard key

% ===========================================================================
\begin{document}

\begin{titlepage}
\centering
\vspace*{3cm}
{\Huge\bfseries\color{accent} Product Intelligence Factory\par}
\vspace{0.6cm}
{\LARGE Demo Guide\par}
\vspace{2cm}
{\large How to run the demonstration, end to end,\\
without having seen the application before.\par}
\vspace{3cm}
\begin{minipage}{0.8\textwidth}
\centering\itshape\color{quiet}
Every command and every button name in this guide was checked against the
running application. If something on screen does not match what is written
here, trust the screen and tell the author --- do not improvise in front of
the audience.
\end{minipage}
\vfill
{\color{quiet}\small Version 2 --- the estate release\par}
\end{titlepage}

\tableofcontents
\newpage

% ===========================================================================
\section{Read this first}

\subsection{What you are demonstrating, in one paragraph}

A large retailer sells thousands of products. The information about each one ---
its size, its wattage, its ingredients, its photographs --- arrives from many
different computer systems, and those systems disagree, send things late, and
sometimes send things that are simply wrong. This application watches all of
them, works out when something has gone wrong, works out everywhere that error
has spread to, fixes it, and stops anything reaching a customer until a human
being has approved the fix.

\subsection{Your job in the room}

You are driving. You are not expected to answer deep technical questions, and
you should not try to. Your job is to:

\begin{enumerate}[leftmargin=*]
  \item Get the application running \emph{before} anybody walks in.
  \item Click through the five acts in Section~\ref{sec:script}, in order.
  \item Say the short lines in the blue boxes.
  \item Write down any question you cannot answer, and say you will find out.
\end{enumerate}

\begin{say}[The single most useful sentence in this guide]
``That is a good question --- let me write it down and get you a proper
answer.''\\[4pt]
Saying this costs you nothing. Guessing costs you the room.
\end{say}

\subsection{How long it takes}

\begin{longtable}{@{}p{5.2cm}p{3cm}p{7cm}@{}}
\toprule
\textbf{Stage} & \textbf{Time} & \textbf{When to do it} \\
\midrule
\endhead
First-time setup & 20--30 minutes & The day before. Not on the day. \\
Pre-flight check & 5 minutes & 30 minutes before the meeting. \\
The demo itself & 12--15 minutes & With questions, allow 25. \\
\bottomrule
\end{longtable}

\begin{careful}[Do the setup the day before]
The first-time setup downloads things from the internet. Meeting-room wifi is
exactly where that goes wrong. Once it is done once, the demo runs without a
network.
\end{careful}

% ===========================================================================
\section{First-time setup (do this the day before)}

You need a terminal window. On Windows, press \key{Start}, type
\texttt{powershell}, and press \key{Enter}.

\subsection{Step 1 --- Go to the project folder}

\begin{cmd}
cd C:\textbackslash Users\textbackslash Sauris\textbackslash Downloads\textbackslash productonboarding-new
\end{cmd}

Every command in this guide is typed in this folder. If you open a new
terminal later, run this line again first.

\subsection{Step 2 --- Check you are on the right version}

\begin{cmd}
git branch --show-current
\end{cmd}

It must print \texttt{V2}. If it prints anything else:

\begin{cmd}
git checkout V2
\end{cmd}

\subsection{Step 3 --- Install what the application needs}

\begin{cmd}
pip install -r requirements.txt
\end{cmd}

This takes a few minutes and prints a great deal of text. That is normal. You
are looking for it to finish without the word \texttt{ERROR} at the end.

\subsection{Step 4 --- Create the settings file}

\begin{cmd}
copy .env.example .env
\end{cmd}

Then open \texttt{.env} in Notepad. You are looking for two lines near the top:

\begin{cmd}
LITELLM\_BASE\_URL=...\\
LITELLM\_API\_KEY=...
\end{cmd}

These tell the application where to find its AI service. \textbf{Ask whoever
gave you this project what to put here.} If you cannot get them, the demo still
works --- see Section~\ref{sec:nogateway}.

\subsection{Step 5 --- Build the demonstration data}

\begin{cmd}
python scripts/generate\_data.py
\end{cmd}

This invents eight weeks of realistic supplier activity --- product records,
supplier emails, specification documents, marketplace rejections. It prints a
summary ending with a line about the ``pack digest''.

\begin{aside}
The same command always produces exactly the same data. That matters: the demo
you rehearse is the demo the CIO sees, down to the individual email.
\end{aside}

\subsection{Step 6 --- Start it once, to prove it works}

\begin{cmd}
python run.py
\end{cmd}

Wait about fifteen seconds, then open a web browser and go to:

\begin{cmd}
http://127.0.0.1:8000
\end{cmd}

You should see a dark or light page with \ui{Ingest Fabric} at the top and a
diagram of boxes joined by lines. If you do, setup is finished.

To stop it, click on the terminal window and press \key{Ctrl}+\key{C}.

% ===========================================================================
\section{Pre-flight (30 minutes before the meeting)}
\label{sec:preflight}

\subsection{Step 1 --- Start with a clean slate}

\begin{cmd}
cd C:\textbackslash Users\textbackslash Sauris\textbackslash Downloads\textbackslash productonboarding-new\\
python run.py --reset
\end{cmd}

The \texttt{--reset} wipes any previous run so the demo starts from day one.

\begin{careful}[If you see ``file is being used by another process'']
An old copy is still running. Close every terminal window you have open, wait
five seconds, and run the command again. This is the single most common
problem, and it is harmless.
\end{careful}

\subsection{Step 2 --- Open the browser}

Go to \texttt{http://127.0.0.1:8000}.

\subsection{Step 3 --- Check the four things}

\begin{longtable}{@{}p{4.6cm}p{10.4cm}@{}}
\toprule
\textbf{Look at} & \textbf{You want to see} \\
\midrule
\endhead
Top right corner & A green dot labelled \ui{gateway}. Green means the AI
service is reachable. Grey or red means it is not --- go to
Section~\ref{sec:nogateway} and adjust what you say. \\[4pt]
The main diagram & Five columns of boxes: \ui{SYSTEMS}, \ui{SOURCES},
\ui{PRODUCTS}, \ui{VARIANTS}, \ui{CHANNELS}. \\[4pt]
The right-hand panel & The words \ui{The catalog is clean}. If it does not say
this, you did not use \texttt{--reset}. Stop and start again. \\[4pt]
The bar along the bottom & A date reading \ui{2026-09-01} and a counter
reading \ui{0/274}. \\
\bottomrule
\end{longtable}

\subsection{Step 4 --- Make it readable from the back of the room}

Click the \ui{Appearance} button in the top right corner (it looks like a
small sun or moon). Choose \ui{Dark}. Dark projects better and photographs
better.

In the browser, press \key{Ctrl}+\key{+} two or three times to make everything
larger. The layout adapts; nothing will break.

\subsection{Step 5 --- Do a full dry run}

Run through Section~\ref{sec:script} once, alone, start to finish. Then press
\ui{Rewind the tape} at the bottom of the screen to put everything back.

\begin{careful}[Why the dry run matters more than it sounds]
The first time the application asks the AI service a question, it pays for that
answer. It then remembers it. A rehearsed demo is therefore both faster and
identical to the one you rehearsed. Skipping the dry run is the difference
between a twenty-second wait and a five-second one.
\end{careful}

% ===========================================================================
\section{The demo script}
\label{sec:script}

Five acts. Follow them in order. Everything you need to click is named in
\ui{bold}.

% ---------------------------------------------------------------------------
\subsection{Act 1 --- The problem (2 minutes)}

You are on the \ui{Ingest Fabric} screen. Do not click anything yet.

\begin{say}[Say this]
``This is every system that sends us product information. Ten of them. A
supplier portal where somebody types things in by hand. The supplier's own
database. The artwork library. Our finance system. An industry data pool. And
so on.

``Right now the catalogue is clean --- nothing is wrong. Watch what happens
over the next eight weeks of supplier activity.''
\end{say}

\textbf{Do this:} scroll down a little to the panel headed \ui{THE ESTATE}.
Point at two of the cards.

\begin{say}[Say this]
``Each one tells us how well it behaves. The artwork library says
\emph{conforms} --- it is the legal source, so it had better. The data pool
says \emph{3 defect kinds} --- it sends things in its own vocabulary, with the
wrong data types, and it contradicts better sources. That is not us being rude
about a supplier. That is what a data pool integration is actually like.''
\end{say}

% ---------------------------------------------------------------------------
\subsection{Act 2 --- Eight weeks in twenty seconds (2 minutes)}

\textbf{Do this:} at the very bottom of the screen, click the button labelled
\ui{Jump to the correction} (hover over the small buttons; the tooltip tells
you which is which).

Wait about five seconds. The screen fills up.

\begin{say}[Say this]
``That is eight weeks of supplier traffic. Two hundred and seventy-four
events, delivered by ten systems in batches, at irregular times, overlapping
each other --- which is what a real morning looks like.

``The systems are now reporting how much they sent, and how much of it was
defective.''
\end{say}

\textbf{Do this:} point at the \ui{THE ESTATE} panel. It now shows counts. Then
point at the feed of small lines below it.

\begin{say}[Say this]
``Every delivery is logged with which system sent it and what was wrong with
it --- a missing mandatory field, a stale document version, a value that
contradicts a better source. Seven named kinds of defect, and every one of them
is something the system can detect. We did not build a data generator that
produces mess nobody checks.''
\end{say}

\begin{aside}
If a CIO asks whether the timings are real: the \emph{delivery} genuinely is
asynchronous and overlapping. What is fixed is \emph{what} each system sends,
so the demo can be rehearsed. It is worth saying so plainly --- it is a design
decision, not a shortcut.
\end{aside}

% ---------------------------------------------------------------------------
\subsection{Act 3 --- Finding and fixing the error (4 minutes)}

This is the centrepiece. Do not rush it.

\textbf{Do this:} click the blue button in the top right, \ui{Run the
correction loop}.

A large banner appears at the top of the screen and starts narrating.

\begin{say}[Say this, while it runs]
``It is working. And notice what it is showing us --- not a progress bar.
Those are its own findings, appearing as it makes them. \emph{Reading the
document. This one is immaterial. This one is a specification correction, one
value read.} You can watch it think.''
\end{say}

Wait for it to finish. It takes about twenty seconds. The screen then changes
to \ui{Blast Radius} on its own.

\textbf{Do this:} let them read the top of the screen, then scroll down slowly.

\begin{say}[Say this]
``It found the correction, worked out which product it applies to, and then
worked out everywhere the old value had already been used --- every bullet
point, every listing, every channel.

``That is not a guess. Every piece of marketing copy in this system records
which values it was built from, so this is a lookup, not an opinion.''
\end{say}

\textbf{Do this:} scroll to the panel headed \ui{WHO HAS TO BE TOLD}.

\begin{say}[Say this]
``And here it is in the words a buyer actually uses. Not internal codes ---
SKUs, and the systems that have to be told about them. If we dispatched this
now, these channels would take it and this one would not, because the printed
catalogue is inside its freeze window and a printed page cannot be recalled.

``It tells you that \emph{before} you decide, not afterwards.''
\end{say}

% ---------------------------------------------------------------------------
\subsection{Act 4 --- The human gate (3 minutes)}

\textbf{Do this:} click \ui{Review \& Audit} in the left-hand menu. There is a
small number badge on it.

\begin{say}[Say this]
``Nothing has been published. It has stopped and is waiting for a person.''
\end{say}

\textbf{Do this:} point at the pink warning box near the top.

\begin{say}[Say this]
``It is telling us why it will not proceed on its own: this is a regulated
food product, its declarations are legally binding, and there is no automatic
path for this change.''
\end{say}

\textbf{Do this:} scroll down slowly through the green list headed \ui{WHAT
GETS WRITTEN IF YOU APPROVE}.

\begin{say}[Say this]
``Every single thing it would change, listed before you decide, each one
citing the supplier document it came from.''
\end{say}

\textbf{Do this:} point at the diagram on the right-hand side.

\begin{say}[Say this]
``And that is the actual decision path it took --- read out of the running
system, not a picture somebody drew. The highlighted box is where it stopped.''
\end{say}

\textbf{Do this:} scroll to the bottom. Type a name into the \ui{Reviewer} box
(use the CIO's name --- it is a nice touch). Click \ui{Approve and republish}.

\begin{say}[Say this]
``Approved under a name, recorded in an audit trail that cannot be edited.
Now it publishes --- and only now.''
\end{say}

% ---------------------------------------------------------------------------
\subsection{Act 5 --- Is a product ready to sell? (4 minutes)}

\textbf{Do this:} click \ui{Product 360} in the left-hand menu.

\begin{say}[Say this]
``Everything so far was about fixing something already published. This is the
question that comes first: is this product good enough to launch at all?''
\end{say}

\textbf{Do this:} point at the top right, where it says how many are ready and
how many are going back.

\textbf{Do this:} in the list on the left, click \ui{Trail Mix Bar Multipack
6x40g}. It has an orange number badge.

\begin{say}[Say this]
``This one is not ready. And look at what it says --- not `data quality score
72\%'. It says the category requires an ingredient panel photograph, none was
supplied, and the imaging system is the one that has to fix it.

``A number would invite somebody to launch at ninety per cent. A named finding
tells you who to email.''
\end{say}

\textbf{Do this:} scroll down to \ui{THE RECORD}.

\begin{say}[Say this]
``And for every value, which system gave it to us. When two systems disagree,
we can say which one was wrong.''
\end{say}

\textbf{Do this:} click \ui{AeroPure 300} at the top of the list --- a green
one. Then type a name into the small box next to the button, and click
\ui{Open staging page}.

\begin{say}[Say this]
``This one passed, so we can see the page it would become.

``And this box at the top is the interesting part. It is written by AI, and it
is fenced: it can only use a value this product actually has, and a market
document we actually wrote. It shows you both, underneath. If it cannot
support a claim from both, it does not make the claim.''
\end{say}

\begin{careful}[The staging page asks for a name]
This is deliberate --- it is unpublished commercial content, and the system
records who looked at it. If you click the button without typing a name,
nothing happens and you will look stuck. Type a name first.
\end{careful}

% ---------------------------------------------------------------------------
\subsection{Optional Act 6 --- Plug and play (2 minutes)}

Only if you have time and the room is engaged.

\textbf{Do this:} go back to \ui{Ingest Fabric}. Scroll to \ui{THE ESTATE}.
Click \ui{Disconnect} on any system.

\begin{say}[Say this]
``Watch the diagram.''
\end{say}

The system disappears from the \ui{SYSTEMS} column immediately.

\begin{say}[Say this]
``Systems can be added and removed while it is running, and the map follows.
Adding one is pasting an address --- there is a box for it just there. Nothing
about the ten systems is written into the application.''
\end{say}

\textbf{Do this:} click \ui{System Control} in the menu and point at the
\ui{Capabilities} panel.

\begin{say}[Say this]
``And everything this system can do publishes itself at a standard address, so
another organisation's software can discover it. Including what each part is
\emph{not} allowed to do --- none of these can approve anything, and none of
them can publish.''
\end{say}

% ===========================================================================
\section{What each screen is for}

Use this if somebody asks you to go back to something.

\begin{longtable}{@{}p{3.4cm}p{11.6cm}@{}}
\toprule
\textbf{Screen} & \textbf{What it shows} \\
\midrule
\endhead
\ui{Ingest Fabric} & The ten external systems, what each is delivering, and the
map of everything. The clock at the bottom controls time. \\[4pt]
\ui{Product 360} & Search for a product; see everything every system has said
about it; see whether it is fit to launch; preview the page it would become. \\[4pt]
\ui{Blast Radius} & For one correction: what it means, which product it applies
to, and everywhere the old value had already been used. \\[4pt]
\ui{Readings} & When a supplier's wording is ambiguous, the competing
interpretations, each scored against the rules. \\[4pt]
\ui{Review \& Audit} & The decision waiting for a human, everything it would
change, and the permanent record of what was published. \\[4pt]
\ui{System Control} & The clock, the AI service, the search index, and the
directory of everything the system can do. \\
\bottomrule
\end{longtable}

\subsection{The bar along the bottom}

It is on every screen. Left to right: play, step forward one event, jump to the
correction, rewind. \ui{Rewind the tape} is your undo --- it puts the demo back
to the beginning.

\subsection{The search box at the top}

Press \key{Ctrl}+\key{K} anywhere to open it. It can jump to any screen and
search the company's written policies. Useful if you get lost.

% ===========================================================================
\section{If something goes wrong}

\subsection{The page will not load}

Check the terminal window is still running. If it has stopped, run
\texttt{python run.py} again. Give it fifteen seconds before refreshing.

\subsection{``The process cannot access the file''}

An old copy is still running. Close every terminal window, wait five seconds,
start again. Harmless.

\subsection{The gateway dot is not green}
\label{sec:nogateway}

The AI service is unreachable. \textbf{The demo still works} --- everything
still runs, using the system's built-in fallbacks. Three things change:

\begin{itemize}[leftmargin=*]
  \item The readiness screen shows an orange note saying it ran without AI.
  \item The staging page's highlighted box is plainer.
  \item Everything is faster.
\end{itemize}

\begin{say}[Say this, and it is a strength not an excuse]
``The AI service is not reachable from this room, so what you are watching is
the system running \emph{without} it. Every step has a deterministic fallback,
and it tells you when it used one. It never quietly presents a guess as an
answer.''
\end{say}

\subsection{A screen is empty or stuck}

Refresh the browser (\key{F5}). Nothing is lost --- the state lives in the
application, not the page.

\subsection{You clicked something and cannot get back}

Click \ui{Rewind the tape} at the bottom, then \ui{Jump to the correction} to
return to Act 3.

% ===========================================================================
\section{Questions a CIO is likely to ask}

Short answers. If pressed further than this, write the question down.

\subsection{``Is the AI making the decisions?''}

\begin{say}
``No, and deliberately not. The AI reads documents and writes prose. Every
decision --- what applies to which product, what is affected, whether something
can be published --- is ordinary software with a fixed answer. The same input
always gives the same output.''
\end{say}

\subsection{``What if the AI is wrong?''}

\begin{say}
``Two things. Every AI finding has to cite a document you can open, and if it
cannot cite one it is discarded. And nothing publishes without a person
approving it.''
\end{say}

\subsection{``How does it connect to our systems?''}

\begin{say}
``Through an open standard called MCP. Each of the ten systems is a separate
endpoint, and connecting a new one is pasting an address --- no code change.
The same on the publishing side.''
\end{say}

\subsection{``Could a connected system do damage?''}

\begin{say}
``Connecting a system lets us \emph{see} what it offers. It does not let the AI
call any of it. Somebody has to allow each tool individually, and that decision
is recorded in the audit trail.''
\end{say}

\subsection{``Is this locked to one AI vendor?''}

\begin{say}
``No. There is no AI model named anywhere in the application --- it uses
whatever the configured service offers, decided when it starts. There is
actually a test that fails the build if somebody hard-codes a model name.''
\end{say}

\subsection{``How fast is it?''}

\begin{say}
``About twenty seconds for a full correction, from reading the supplier
document to stopping at the approval gate.''
\end{say}

\subsection{``Is it tested?''}

\begin{say}
``446 automated tests. They run with the AI service deliberately switched off,
so what is tested is the behaviour when things are not working.''
\end{say}

% ===========================================================================
\section{What not to claim}

Read this section twice. A CIO will forgive a limitation. They will not forgive
finding one you glossed over.

\begin{longtable}{@{}p{6.4cm}p{8.6cm}@{}}
\toprule
\textbf{Do not say} & \textbf{Say instead} \\
\midrule
\endhead
``It's fully automated.'' & ``It automates everything up to the decision, and
stops there on purpose.'' \\[4pt]
``It connects to any system.'' & ``It connects to anything that speaks MCP.
Systems that don't would need an adapter.'' \\[4pt]
``The AI checks everything.'' & ``Six of the nine checks are rules. Three use
AI, and those have to cite a document.'' \\[4pt]
``It's production ready.'' & ``It's a working prototype with a real test suite.
There's no login yet --- decisions are recorded against a typed name.'' \\[4pt]
``The images are real.'' & ``Product photographs are represented in the record,
not stored. It checks whether one was supplied, not what is in it.'' \\
\bottomrule
\end{longtable}

\subsection{One honest technical caveat, if a technical person is present}

If somebody asks specifically about the search technology, this is worth
knowing and worth saying:

\begin{say}
``The search combines keyword and AI-based matching. The keyword half is
measurably doing the work --- it separates product codes that differ by one
character, which is exactly where this matters. The AI half is not currently
adding measurable value on our test questions, and we know that because we
measured it rather than assumed it. It's written down as an open item.''
\end{say}

Saying this unprompted, in front of a CIO, will do more for your credibility
than any feature in the demo.

% ===========================================================================
\section{One-page cheat sheet}

Tear this page off and keep it next to you.

\vspace{0.6em}
\begin{tcolorbox}[colback=shell,colframe=quiet,boxrule=0.4pt,arc=2pt]
\textbf{Before they arrive}\\[2pt]
\texttt{cd C:\textbackslash Users\textbackslash Sauris\textbackslash Downloads\textbackslash productonboarding-new}\\
\texttt{python run.py --reset}\\
Browser: \texttt{http://127.0.0.1:8000}\\
Check: green \ui{gateway} dot, \ui{The catalog is clean}, counter \ui{0/274}
\end{tcolorbox}

\vspace{0.4em}
\begin{tcolorbox}[colback=white,colframe=accent,boxrule=1pt,arc=2pt]
\textbf{The five acts}\\[4pt]
\textbf{1.} \ui{Ingest Fabric} --- ``ten systems, catalogue is clean'' \\[3pt]
\textbf{2.} \ui{Jump to the correction} (bottom bar) --- ``eight weeks of
traffic'' \\[3pt]
\textbf{3.} \ui{Run the correction loop} (blue, top right) --- wait 20s ---
``watch it think'', then scroll to \ui{WHO HAS TO BE TOLD} \\[3pt]
\textbf{4.} \ui{Review \& Audit} --- ``it stopped for a human'' --- scroll,
type a name, \ui{Approve and republish} \\[3pt]
\textbf{5.} \ui{Product 360} --- click the orange one, ``named finding, not a
score'' --- click a green one, type a name, \ui{Open staging page}
\end{tcolorbox}

\vspace{0.4em}
\begin{tcolorbox}[colback=warnbg,colframe=warnline,boxrule=0.4pt,arc=2pt]
\textbf{If it breaks}\\[2pt]
Page dead $\rightarrow$ check the terminal, \texttt{python run.py} again\\
``file in use'' $\rightarrow$ close all terminals, retry\\
Gateway not green $\rightarrow$ carry on, say it is running without AI\\
Lost $\rightarrow$ \ui{Rewind the tape}, then \ui{Jump to the correction}\\
Stuck on staging page $\rightarrow$ you forgot to type a name
\end{tcolorbox}

\vspace{0.4em}
\begin{tcolorbox}[colback=saybg,colframe=accent,boxrule=0pt,leftrule=3pt,arc=0pt]
\textbf{The three lines that matter most}\\[4pt]
``The AI reads. The software decides. A person approves.''\\[3pt]
``Every AI finding cites a document you can open.''\\[3pt]
``That's a good question --- let me write it down and get you a proper
answer.''
\end{tcolorbox}

\end{document}
